% Reference-motion synthesis for microspine wall climbing (method section draft)
% Build: xelatex reference_synthesis_method.tex   (pdflatex also works)
%
% Provenance note (not part of the manuscript): Secs. 3, 5, 6 and the
% weighted variant of Sec. 4.6 describe what scripts/gen_climb_clip.py,
% scripts/reference/gen_mount_clip.py, scripts/validate_climb_reference.py and the
% climb_wall MDP implement today (2026-09-04). The lexicographic solver of
% Secs. 4.1-4.5 is the target formulation agreed in conversation; its
% implementation is planned (see docs/plans). Numbers marked [TBD] are
% placeholders for the ablation.
\documentclass[11pt]{article}
\usepackage[a4paper,margin=2.2cm]{geometry}
\usepackage{amsmath,amssymb,bm}
\usepackage{booktabs}
\usepackage{array}
\usepackage{xcolor}
\usepackage{enumitem}
\usepackage[hidelinks]{hyperref}
\setlength{\parskip}{3pt}

\newcommand{\R}{\mathbb{R}}
\newcommand{\SO}{\mathrm{SO}(3)}
\newcommand{\SE}{\mathrm{SE}(3)}
\newcommand{\nhat}{\hat{\bm n}}
\newcommand{\that}{\hat{\bm t}}
\newcommand{\lhat}{\hat{\bm l}}
\newcommand{\zhat}{\hat{\bm z}}
\newcommand{\dq}{\delta\bm q}
\newcommand{\vect}[1]{\bm{#1}}
\newcommand{\tbd}{{\color{red}[TBD]}}
\DeclareMathOperator*{\argmin}{arg\,min}
\DeclareMathOperator{\diag}{diag}
\DeclareMathOperator{\atan}{atan2}

\title{Contact-Scheduled Kinematic Reference Synthesis\\for Microspine Wall Climbing without Motion Capture}
\author{Method section draft}
\date{2026-09-04}

\begin{document}
\maketitle

\section{Problem statement}

\subsection{Setting and notation}

We consider a humanoid (PHYBot~C2) climbing a planar wall inclined at angle
$\alpha$ from the floor ($\alpha=90^\circ$: vertical). The robot has
$n_j=21$ actuated joints $\bm\theta\in\R^{n_j}$ (6 per leg, 4 per wristless
arm, 1 waist yaw) and a floating base $(\bm p_b,R_b)\in\SE$. Four microspine
pads $h\in\mathcal H=\{\mathrm{LH},\mathrm{RH},\mathrm{LF},\mathrm{RF}\}$
are mounted on the hands and shins. Each pad sits on a passive three-hinge
gimbal with angles $\bm\gamma_h\in\R^3$, mechanical stops
$|\gamma_{h,k}|\le\gamma_{\mathrm{stop}}=3^\circ$, and a weak return spring
($0.7\,\mathrm{N\,m/rad}$). The feet additionally carry a rigid toe tip that
may press, but not adhere to, the wall.

Positions are expressed in the canonical wall frame: the wall surface is the
plane $x=x_w$, $\zhat$ points up the wall, and $\nhat=-\hat{\bm x}$ is the
outward wall normal (pointing at the robot). A tilted wall is obtained by
rotating this frame about the wall base line; gravity stays physically
vertical. On an engaged pad the \emph{engage frame}
$(\that,\lhat,\nhat)$ has $\that$ along the rake direction (down-wall),
$\lhat=\nhat\times\that$ lateral.

A microspine pad holds a load $\bm F=(f_t,f_l,f_n)$ only inside the closed,
direction-dependent limit surface of Wang, Jiang and Cutkosky~(IJRR~2017):
capacity peaks along the rake shear $f_t>0$, tension $f_n>0$ is admissible
only alongside shear, and against-rake and pure normal pull-off hold almost
nothing. A pad seats only after a small down-wall stroke under compression
(the \emph{rake}) and recruits its full capacity over
$\tau_{\mathrm{ramp}}\approx0.25\,\mathrm s$ of loaded dwell.

A reference motion is a tuple
\begin{equation}
  \mathcal M=\bigl(\bm q(t),\;\bm\sigma(t),\;\{A_h(t)\}_{h\in\mathcal H},\;\bm\beta(t)\bigr),
  \qquad t\in[0,T],
\end{equation}
with configuration $\bm q=(\bm p_b,R_b,\bm\theta)$, contact schedule
$\bm\sigma(t)\in\{0,1\}^{4}$ ($\sigma_h=1$: pad $h$ latched), anchor poses
$A_h=(\bm a_h,R^a_h)\in\SE$ (the weld target while latched, the landing
target while swinging), and a toe-brace schedule $\bm\beta(t)\in\{0,1\}^2$.
The reference is consumed by reinforcement learning as an imitation target
for \emph{shape and schedule only}: no reward tracks the base position in the
world, and force allocation and balance are left to the learned policy.

\subsection{Why motion capture does not transfer}

Imitation pipelines for humanoids obtain $\bm q(t)$ from human motion
capture followed by kinematic retargeting. Two obstacles make this route
unavailable here.

\paragraph{Data.} Adhesive wall climbing with palm and shin pads is not a
human activity; no capture corpus exists, and staged approximations
(bouldering, ladder climbing) exercise fingers, wrists and toes the robot
does not have and load the holds through friction geometry the pads cannot
reproduce.

\paragraph{Embodiment optimality.} Even with data, the human optimum is not
the robot optimum. The feasible set of the robot is dictated by (i)~a
4-DoF wristless arm that cannot satisfy an arbitrary 6-DoF hand pose,
(ii)~a pad that must stay flat on the wall and may rotate only within the
$3^\circ$ gimbal travel once latched, (iii)~an adhesion envelope that is
essentially one-sided: the pad must be loaded along its rake direction and
must never be dragged across the surface while unloaded, and (iv)~a
recruitment time that forbids fast regrips. Retargeting projects a human
trajectory onto the robot's kinematic chain, but not onto this
\emph{contact-feasibility} set; the projected motion drags latched pads,
loads them against the rake, or transfers weight onto immature anchors.
In our experience, imitating references that lack an explicit schedule and
the physics gates of Sec.~\ref{sec:gates} produced degenerate gaits: policies
welded two limbs permanently and dragged the others (Sec.~\ref{sec:ablation}).

\subsection{Requirements}

A usable reference must satisfy:
\begin{enumerate}[label=(R\arabic*),leftmargin=2.6em]
  \item \textbf{Explicit schedule.} $\bm\sigma$ and $\{A_h\}$ are part of the
        reference, so that the policy can be rewarded for latching the
        right pad at the right place and time.
  \item \textbf{Kinematic closure.} While $\sigma_h=1$ the pad is
        stationary at $A_h$ (a latched pad physically cannot move), the
        limb may rotate relative to it only within the gimbal travel, all
        joints stay strictly inside their limits, and released pads and toes
        do not penetrate the wall.
  \item \textbf{Static feasibility.} At every frame there exists a force
        allocation on the scheduled contacts that balances gravity inside
        the limit surface and the actuator torque limits.
  \item \textbf{Timing feasibility.} Dwell times respect the engage pipeline
        and the recruitment ramp.
  \item \textbf{Periodicity.} One cycle ends exactly at its start
        translated by the step height $h$ up the wall, so the clip loops.
  \item \textbf{Determinism and low cost.} Regeneration takes minutes and
        is independent of runtime physics parameters, so that the reference
        can sit inside a design loop.
\end{enumerate}

\section{Method overview}

The reference is synthesized in three stages (Fig.~\tbd).
\emph{Choreography} (Sec.~\ref{sec:choreo}) fixes the contact schedule and a
small set of task-space intents from the robot's own stance.
\emph{Lexicographic constrained inverse kinematics}
(Sec.~\ref{sec:ik}) turns the intents into joint trajectories frame by
frame, with the latched contacts and the mechanical limits as hard
constraints and every intent in their null space. \emph{Physics gates}
(Sec.~\ref{sec:gates}) certify static and timing feasibility and close the
loop on the choreography parameters. Section~\ref{sec:rl} states how the
policy consumes the result.

\section{Choreography: contact schedule and task-space intents}
\label{sec:choreo}

\paragraph{Stance and anchors.} The cycle starts from the climb stance
$\bm q^0$ (all four pads latched, arms overhead, legs splayed), obtained
once by an inverse-kinematics stance search that pins the four pad frames
flat on the wall plane at a chosen hand height and base standoff and
minimizes posture cost. All anchors are derived from it by forward
kinematics,
\begin{equation}
  A^0_h=\mathrm{FK}_h(\bm q^0),\qquad
  A^+_h=\bigl(\bm a^0_h+h\,\zhat,\;R^{a,0}_h\bigr),\qquad h=0.03\,\mathrm m,
\end{equation}
so that no world coordinate is authored by hand except the wall plane.

\paragraph{Schedule.} One cycle is a four-beat crawl,
\[
\text{hold}\to\text{LF swing}\to\text{hold}\to\text{RF swing}\to\text{hold}
\to\text{LH swing}\to\text{hold}\to\text{RH swing}\to\text{hold},
\]
subject to the rules that encode the pad physics:
\begin{enumerate}[label=(S\arabic*),leftmargin=2.6em]
  \item at most one pad is released at any time, so $\sum_h\sigma_h\ge3$;
  \item the support set always contains at least one hand and one foot,
        preserving the upper-tension/lower-compression couple that balances
        the wall-normal gravity moment;
  \item every hand swing starts with a short all-four load-transfer dwell,
        and every dwell before a release lasts at least $\tau_{\mathrm{ramp}}$;
  \item during a foot swing the toe of the released foot braces the wall as a
        compression contact ($\beta_h=1$), so the unloaded spine pad never
        slides;
  \item every landing ends with a rake: the pad first touches the wall
        $d_{\mathrm{rake}}$ above its anchor and is dragged straight down the
        surface onto it before $\sigma_h$ flips to~1.
\end{enumerate}
Segment durations $N_k$ (frames at $f_s=50$\,Hz) are choreography parameters;
Table~\ref{tab:segments} lists the shipped cycle.

\begin{table}[h]
\centering\small
\begin{tabular}{@{}lccl@{}}
\toprule
Segment & Frames & Duration [s] & Released pad / brace \\
\midrule
hold\_start & 15 & 0.30 & none (left toe plants) \\
left\_foot\_swing & 38 & 0.76 & LF, left toe braces \\
foot\_transfer\_hold & 15 & 0.30 & none (retract left toe, plant right) \\
right\_foot\_swing & 38 & 0.76 & RF, right toe braces \\
seat\_hold & 15 & 0.30 & none (retract right toe) \\
left\_hand\_swing & 45 & 0.90 & LH after 0.1\,s transfer dwell \\
regrip\_hold & 15 & 0.30 & none \\
right\_hand\_swing & 45 & 0.90 & RH after 0.1\,s transfer dwell \\
\midrule
cycle & 226 & 4.52 & $h=30$\,mm, $6.6$\,mm/s \\
\bottomrule
\end{tabular}
\caption{Contact schedule of the shipped reference cycle
(\texttt{climb\_cycle\_moment\_safe\_paced}).}
\label{tab:segments}
\end{table}

\paragraph{Time parameterization.} Within a segment of $N$ frames the
progress variable is
\begin{equation}
  s(t)=S\!\left(\tfrac{t}{N-1}\right),\qquad S(x)=6x^5-15x^4+10x^3,
\end{equation}
which has zero velocity and acceleration at both ends.

\paragraph{Swing path.} For the released pad with start pose
$(\bm p^0,R^0)=\mathrm{FK}_h(\bm q)$ at release and target anchor
$(\bm p^+,R^a)$, the last $N_r$ frames are the rake and the first
$N_a=N-N_r$ the approach arc,
\begin{align}
  \bm c &= \tfrac12(\bm p^0+\bm p^+)+b\,\nhat+\ell\,\hat{\bm y}_{\mathrm{side}},
  \qquad \bm r=d_{\mathrm{rake}}\,\zhat,\\
  \bm p^*(t) &=
  \begin{cases}
    (1-s_a)^2\,\bm p^0+2(1-s_a)s_a\,\bm c+s_a^2\,(\bm p^++\bm r), &
    t<N_a,\quad s_a=S\!\bigl(\min(\tfrac{t}{N_a-1},1)\bigr),\\[2pt]
    \bm p^++\bigl(1-S(\tfrac{t-N_a+1}{N_r})\bigr)\,\bm r, & t\ge N_a,
  \end{cases}
\end{align}
with off-wall bulge $b$ (hands $0.04$\,m, feet $0.03$\,m), lateral offset
$\ell$ (default~0) and $d_{\mathrm{rake}}=2$\,mm. The orientation target of a
foot is the geodesic interpolation
$R^*(s)=\exp\!\bigl(s\log(R^aR^{0\top})\bigr)R^0$; a hand holds $R^0$ and
its orientation weight ramps toward landing (Sec.~\ref{sec:soft}).

\paragraph{Toe brace.} With toe tip position $\bm p_{\mathrm{toe}}$ and radius
$r_{\mathrm{toe}}$ the brace target is
$x^*_{\mathrm{toe}}(t)=(1-\pi(t))\,x_{\mathrm{plant}}+\pi(t)\,(x_w-r_{\mathrm{toe}})$,
where the press profile $\pi(t)$ ramps $0\!\to\!1$ over $0.06$\,s, holds 1,
and returns to 0 during the rake; $y,z$ ride freely.

\paragraph{Base and posture intents.} During a swing the base target and the
posture reference blend from their values at the segment start,
\begin{equation}
  \bm p^*_b(s)=(1-s)\,\bm p_b(t_0)+s\,\bm p^{\mathrm{end}}_b,\qquad
  \bm\theta^*(s)=(1-s)\,\bm\theta(t_0)+s\,\bm\theta^0,
\end{equation}
with $\bm p^{\mathrm{end}}_b=\bm p^0_b-\delta_{\mathrm{press}}\nhat
+\delta_{\mathrm{lat}}\hat{\bm y}_{\mathrm{side}}+\tfrac{k+1}{2}h\,\zhat$ for
the $k$-th foot swing ($\delta_{\mathrm{press}}=0.03$\,m toward the wall,
$\delta_{\mathrm{lat}}=0.025$\,m toward the supporting side). These are
intents, not constraints: the solver realizes them only as far as the latched
contacts allow.

\section{Lexicographic constrained inverse kinematics}
\label{sec:ik}

Each frame is solved by iterating a two-level quadratic program in the
tangent increment of the configuration until convergence. The design
principle is that \emph{physical facts are hard and intents are soft}: a
latched pad does not move, a hinge does not pass its stop, a joint does not
leave its range; everything the choreographer merely \emph{wants} is realized
in the null space of those facts.

\subsection{Decision variables and hard bounds (level 0)}

The step is
\begin{equation}
  \dq=(\delta\bm p_b,\;\delta\bm\phi_b,\;\delta\bm\theta,\;\delta\bm\gamma)\in\R^{n},
  \qquad n=6+n_j+3|\mathcal H|=39,
\end{equation}
where $\delta\bm\phi_b$ is the base rotation vector and $\delta\bm\gamma$
stacks the gimbal angles of all four pads. Bounds encode the mechanical
limits and a trust region:
\begin{align}
  \bm\theta_{\mathrm{lo}}+\bm m\;\le\;\bm\theta+\delta\bm\theta\;&\le\;\bm\theta_{\mathrm{hi}}-\bm m,
  & \bm m&=\mu\,(\bm\theta_{\mathrm{hi}}-\bm\theta_{\mathrm{lo}}),\ \mu=0.01,\label{eq:jl}\\
  -\gamma_{\max}\le\gamma_{h,k}+\delta\gamma_{h,k}&\le\gamma_{\max}
  & &\text{if }\sigma_h=1,\quad \gamma_{\max}<\gamma_{\mathrm{stop}},\label{eq:gim}\\
  \gamma_{h,k}+\delta\gamma_{h,k}&=0 & &\text{if }\sigma_h=0,\label{eq:gim0}\\
  \|\dq\|_\infty&\le\Delta. & &\label{eq:tr}
\end{align}
The joint margin $\mu$ matches the $0.98$ soft band the downstream policy
clamps its targets to, so the reference never asks the policy to fight its
own clamp. A latched pad may use its gimbal travel up to $\gamma_{\max}$
(default $2^\circ$), leaving $\gamma_{\mathrm{stop}}-\gamma_{\max}$ of travel
for the policy's own tracking error. A released pad's gimbal is pinned at
zero: in the plant it is passive and unloaded, and the solver must not use
it to fake a landing orientation. All of~\eqref{eq:jl}--\eqref{eq:tr} are
box constraints $\bm l\le\dq\le\bm u$.

\subsection{Hard task-space constraints (level 1)}

For every latched pad $h$ ($\sigma_h=1$) the pad frame, i.e.\ the body
\emph{below} the gimbal, is welded to its anchor. With
$(\bm p_h,R_h)=\mathrm{FK}_h(\bm q)$, site Jacobians $J^p_h,J^r_h\in\R^{3\times n}$
(including the gimbal columns), and the world-frame small-angle error
\begin{equation}
  \bm e^p_h=\bm a_h-\bm p_h,\qquad
  \bm e^r_h=\tfrac12\sum_{k=1}^{3}\bm r_{h,k}\times\bm r^a_{h,k}
  \quad(\bm r_{h,k},\bm r^a_{h,k}\text{: columns of }R_h,R^a_h),
\end{equation}
the weld rows are
\begin{equation}
  J_1\,\dq=\bm e_1,\qquad
  J_1=\begin{bmatrix}J^p_{h_1}\\ J^r_{h_1}\\ \vdots\end{bmatrix},\quad
  \bm e_1=\begin{bmatrix}\bm e^p_{h_1}\\ \bm e^r_{h_1}\\ \vdots\end{bmatrix}
  \in\R^{6n_c},
  \label{eq:weld}
\end{equation}
$n_c=\sum_h\sigma_h\in\{3,4\}$. Wall clearance of released pads and toes is
a linear inequality on the same level: for every contact sphere or toe tip
$c$ of a released pad with center $\bm p_c$, radius $r_c$ and Jacobian $J_c$,
\begin{equation}
  \hat{\bm x}^\top\bigl(\bm p_c+J_c\,\dq\bigr)\le x_w-r_c
  \quad\Longleftrightarrow\quad G\,\dq\le\bm g .
  \label{eq:clear}
\end{equation}

\subsection{Soft objectives in the null space (level 2)}
\label{sec:soft}

The intents form a weighted least-squares cost,
\begin{equation}
\begin{aligned}
  f_2(\dq)={}&\tfrac12\!\!\sum_{i\in\mathcal S}\!\bigl\|W_i(J_i\dq-\bm e_i)\bigr\|^2
  +\tfrac12\bigl\|W_\theta(\delta\bm\theta-(\bm\theta^*-\bm\theta))\bigr\|^2\\
  &+\tfrac12\bigl\|W_{\mathrm{lim}}(\delta\bm\theta-(\bm\theta_{\mathrm{bnd}}-\bm\theta))\bigr\|^2
  +\tfrac12 w_\gamma^2\|\bm\gamma+\delta\bm\gamma\|^2
  +\tfrac12\lambda^2\|\dq\|^2,
\end{aligned}
\label{eq:f2}
\end{equation}
where $\mathcal S$ collects the swing-pad pose task, the toe-brace task and
the base task; $W_\theta=\diag(w_\theta)$ pulls toward the posture reference;
$W_{\mathrm{lim}}$ is nonzero only for joints inside a margin band of width
$\mu_s(\theta_{\mathrm{hi}}-\theta_{\mathrm{lo}})$, $\mu_s=0.10$, next to
either stop and pulls them back to the band boundary $\theta_{\mathrm{bnd}}$;
$w_\gamma$ prefers the arm joints over the gimbal; and $\lambda^2=0.5$ is the
Levenberg damping. Table~\ref{tab:weights} lists the weights. The swing hand's
weights fade in over the first 30\,\% of the arc and grow again toward
landing,
$w=w_{\min}+(w_{\max}-w_{\min})\max\!\bigl(1-\min(s_a/0.3,1),\,s_a^2\bigr)$,
which spreads the re-equilibration caused by releasing a weld over the whole
swing instead of a one-frame pop.

\begin{table}[h]
\centering\small
\begin{tabular}{@{}llcc@{}}
\toprule
Task & Rows & $w_{\mathrm{pos}}$ & $w_{\mathrm{rot}}$ \\
\midrule
latched pad (weld) & level 1, hard & -- & -- \\
swing foot pad & 6 & 1000 & 100 \\
swing hand pad & 6 & $500\to2000$ & $50\to500$ \\
toe brace & 3 (position only) & $(300,\,30,\,30)$ & -- \\
base intent & 6 & $30$--$100$ & $2$--$50$ \\
posture reference $w_\theta$ & $n_j$ & $0.3$--$0.6$ (RH swing 10) & -- \\
limit-band spring $w_{\mathrm{lim}}$ & active joints & 3 & -- \\
gimbal return $w_\gamma$ & 12 & 3 & -- \\
damping $\lambda$ & $n$ & $\sqrt{0.5}$ & -- \\
\bottomrule
\end{tabular}
\caption{Task hierarchy and level-2 weights. Latched pads are hard
constraints and carry no weight.}
\label{tab:weights}
\end{table}

\subsection{Two-stage solution, convergence and feasibility verdict}

Each iteration solves two convex QPs. The first realizes the hard tasks as
well as the bounds allow,
\begin{equation}
  \dq_1=\argmin_{\dq}\ \tfrac12\|J_1\dq-\bm e_1\|^2
  \quad\text{s.t.}\quad \bm l\le\dq\le\bm u,\ \ G\dq\le\bm g,
  \qquad \bm r_1=J_1\dq_1-\bm e_1 .
  \label{eq:qp1}
\end{equation}
The second minimizes the intents without degrading that optimum:
\begin{equation}
  \dq_2=\argmin_{\dq}\ f_2(\dq)
  \quad\text{s.t.}\quad U_r^\top J_1\dq=U_r^\top J_1\dq_1,\ \
  \bm l\le\dq\le\bm u,\ \ G\dq\le\bm g,
  \label{eq:qp2}
\end{equation}
where $U_r$ holds the left singular vectors of $J_1$ with
$\sigma_i>\varepsilon\sigma_{\max}$, so that rank-deficient weld stacks
(e.g.\ four latched pads on a closed chain) yield a consistent equality set.
The configuration is updated on the manifold,
$\bm q\leftarrow\bm q\oplus\dq_2$ (base rotation by the exponential map,
joints and gimbals additively), and the iteration stops when
\begin{equation}
  \|\bm e^p_h\|<\varepsilon_p=1.5\times10^{-4}\,\mathrm m,\quad
  \|\bm e^r_h\|<\varepsilon_r=2\times10^{-3}\,\mathrm{rad}\ \ \forall h:\sigma_h=1,\quad
  \|\dq_2\|_\infty<\varepsilon_q=10^{-4},
\end{equation}
or after $K_{\max}=100$ iterations ($\Delta=0.12$; the first frames after a
release are throttled to $\Delta=0.003$ and two iterations).

\paragraph{Lexicographic optimality.} By construction $\dq_2$ attains the
level-1 residual $\bm r_1$ and, among all steps that do, minimizes $f_2$.
With equality constraints only this reduces to the classical null-space
projection $\dq_2=\dq_1+N\bm z^*$, $N=I-V_rV_r^\top$; the QP form is what
admits the box and clearance inequalities at the same priority as the welds.

\paragraph{Feasibility verdict.} If at termination
$\|\bm r_1\|$ exceeds the tolerances, the frame is reported infeasible
together with the active bounds (which joint or gimbal sits at its limit).
The solver never resolves such a conflict by moving a latched pad or by
overshooting a stop; it surfaces the conflict to the choreographer. This is
the mechanism that turns choreography design into a search over a small
parameter vector with a binary kinematic verdict per frame
(Sec.~\ref{sec:loop}).

\subsection{Equivalence: bounded gimbal $\equiv$ box-relaxed rotational weld}

Let $\bm o_h$ be the gimbal pivot (the centroid of the pad's contact points)
and $R^{L}_h$ the orientation of the limb link above the gimbal. Since the pad
is rigidly attached to the wall while latched and the limb is related to it
by the gimbal rotation about $\bm o_h$, constraints
\eqref{eq:weld}--\eqref{eq:gim} are equivalent to
\begin{equation}
  \bm o_h(\bm q)=\bm o^a_h,\qquad
  \bigl|\log\!\bigl(R^{L\top}_h(\bm q)\,R^a_h\bigr)_k\bigr|\le\gamma_{\max},\quad k=1,2,3,
\end{equation}
i.e.\ a 3-DoF positional weld of the limb at the pivot plus a componentwise
box on its orientation error. The bounded-variable form is preferred in the
implementation because it reuses the plant's own Jacobians and pivot geometry
without a small-angle approximation, and because the recorded
$\bm\gamma_h(t)$ directly predicts how much of the physical travel the
reference consumes.

\subsection{Weighted single-level variant (baseline)}
\label{sec:weighted}

The baseline replaces the hierarchy by one weighted damped least-squares
step,
\begin{equation}
  \dq=\Bigl(\sum_i J_i^\top W_i^2J_i+D+\lambda^2 I\Bigr)^{-1}
      \Bigl(\sum_i J_i^\top W_i^2\bm e_i+\bm d\Bigr),
\end{equation}
where the weld rows enter with $W=\diag(1000,1000,1000,500,500,500)$, the
gimbals are locked at zero, and $(D,\bm d)$ collect the posture and
limit-band springs; joint limits are enforced by clipping after the step.
Because weights trade off, the weld leaks: the latched pad drifts by an
amount that grows with the intent weights (measured up to $3.4$\,mm inside
the swings of the shipped cycle and $1.6$\,mm on the mount prologue under a
$2\times$ weld weight; re-anchoring at every segment boundary hides the
accumulated drift from the dwells), and an over-constrained arm resolves a
rotational demand by translating the pad off its anchor rather than by
reporting infeasibility. Section~\ref{sec:ablation}
compares the two.

\subsection{Cycle closure, tiling and velocities}

The last swing returns the robot to the stance; residual solver error is
removed by a closure blend over the final $n_c=\lceil0.4N_{\mathrm{last}}\rceil$
frames,
\begin{equation}
  \bm\Delta=(\bm q^0\oplus h\zhat)\ominus\bm q_{T},\qquad
  \bm q_{T-n_c+i}\leftarrow\bm q_{T-n_c+i}\oplus S\!\left(\tfrac{i+1}{n_c}\right)\bm\Delta,
  \quad i=0,\dots,n_c-1,
\end{equation}
so that $\bm q(T)=\bm q^0\oplus h\zhat$ exactly: same joint angles, base
translated by one step. Cycles tile as $\bm q^{(k)}=\bm q\oplus kh\zhat$ with
anchors shifted alike. Velocities are forward differences on the loop,
$\dot{\bm\theta}_t=f_s(\bm\theta_{t+1}-\bm\theta_t)$ and
$\bm\omega_t=f_s\log(q_{t+1}\otimes q_t^{-1})$, where frame $T$ is frame~0
shifted, so the wrap velocity is the true loop velocity.

\section{Physics gates}
\label{sec:gates}

The generator is purely kinematic. Two audits certify that the result is
physically holdable; both are cheap, deterministic, and independent of any
learned policy and of the simulator's constraint stiffness.

\subsection{Static force-existence audit on the limit surface}

At each frame the robot is posed at $\bm q(t)$ (gimbals at their recorded
angles) and we ask whether \emph{some} allocation of contact forces balances
gravity within the pad envelope and the actuator limits. Unknowns are the
pad-on-wall forces $\bm F_h$ of latched pads, unilateral wall presses
$\bm C_j$ (raking pad faces, braced toes), floor forces $\bm P_k$ (mount
prologues only) and a slack $s$:
\begin{equation}
\begin{aligned}
  \min_{\{\bm F_h\},\{\bm C_j\},\{\bm P_k\},s}\quad & s+\epsilon\|\bm z\|^2\\
  \text{s.t.}\quad
  & \textstyle\sum_h\bm F_h+\sum_j\bm C_j+\sum_k\bm P_k=m\,\bm g,\\
  & \textstyle\sum_h(\bm r_h-\bm r_{\mathrm{com}})\times\bm F_h
    +\sum_j(\bm r_j-\bm r_{\mathrm{com}})\times\bm C_j
    +\sum_k(\bm r_k-\bm r_{\mathrm{com}})\times\bm P_k=\bm 0,\\
  & u_h(\bm F_h)\le s\quad\forall h:\sigma_h=1,\\
  & \bigl|\bm\tau_0+\textstyle\sum_h J_h^\top\bm F_h+\sum_j J_j^\top\bm C_j+\sum_k J_k^\top\bm P_k\bigr|
    \le\bm\tau_{\max},\\
  & c_{j,n}\ge0,\ \|\bm c_{j,t}\|\le\mu_w c_{j,n};\qquad
    P_{k,z}\le0,\ \|\bm P_{k,xy}\|\le\mu_g(-P_{k,z}),
\end{aligned}
\label{eq:audit}
\end{equation}
where $\bm\tau_0$ is the gravity torque at zero velocity minus passive joint
torques, $\bm\tau_{\max}$ the deployment torque limits, and the pad
utilization is
\begin{equation}
  u_h(\bm F)=\frac{\sqrt{f_t^2+f_l^2+\max(f_n,0)^2}}
                  {F_0\,\max\!\bigl(R(\vartheta,\varphi),\,r_{\mathrm{floor}}\bigr)},
  \qquad
  \vartheta=\atan(|f_l|,f_t),\quad
  \varphi=\atan\!\bigl(\max(f_n,0),\sqrt{f_t^2+f_l^2}\bigr),
  \label{eq:util}
\end{equation}
with $R(\vartheta,\varphi)\in[0,1]$ the normalized Wang limit surface
(tabulated, bilinearly interpolated), $F_0$ the pure-rake capacity
($400$\,N, the minimum of the training-time randomization range) and
compression $f_n<0$ carried by ordinary contact. Local pad moments are zero
by construction (point forces at the pad centers), so a passing frame cannot
be relying on the weld as a flange. The audit \emph{passes} iff $s^\star\le1$
at every frame; the shipped cycle reports $\max_t s^\star=0.877$. Margins are
built in: $\mu_w=0.5$ (simulated $0.8$), $\mu_g=0.6$ (simulated $1.0$).

\subsection{Schedule-timing audit}

Let $\tau_{\mathrm{cd}}=0.20$\,s be the re-engage cooldown,
$\tau_{\mathrm{db}}=0.06$\,s the press debounce and
$\tau_{\mathrm{mv}}=0.25$\,s a floor for the actual pad travel. Over the
circular schedule, every OFF window $W$ of every pad must satisfy
\begin{equation}
  |W|\ge\tau_{\mathrm{cd}}+\tau_{\mathrm{db}}+\tau_{\mathrm{mv}},
\end{equation}
and for every regrip of pad $h$ at $t_r$ and the next release $t_f$ of any
other pad,
\begin{equation}
  t_f-t_r\ge\tau_{\mathrm{ramp}}=0.25\,\mathrm s
  \qquad(\text{warning if }t_f-t_r<\tau_{\mathrm{ramp}}+0.3\,|W|),
\end{equation}
because a freshly latched pad holds only a fraction of its capacity until it
has recruited under load. A cycle that violated this by $0.05$\,s
(\texttt{fast} timing, $0.2$\,s seat holds) trained policies that abandoned
the immature foot; restoring the dwells (Table~\ref{tab:segments}) fixed it
without any change of pose.

\subsection{Design loop}
\label{sec:loop}

The choreography is a vector $\bm\psi$ of about fifteen scalars (step height,
bulges, press and lateral shifts, rake length, brace timing) plus the segment
durations. Generation with the IK verdict, the static audit and the timing
audit takes minutes and is bitwise reproducible, so $\bm\psi$ is tuned in a
tight loop and the final recipe is stored with the clip. Because the
generator does not read any runtime physics parameter, recalibrating the
simulator's pad model does not invalidate the reference: after the
$4^\circ$ suspension recalibration the regenerated clip was identical to the
previous one.

\section{Consuming the reference in reinforcement learning}
\label{sec:rl}

The clip is resampled to the policy step $\Delta t$ and indexed by a phase
$k\in\{0,\dots,C-1\}$ that advances every step and wraps with the base
offset increased by $h$. Episodes start by reference-state initialization: a
uniformly sampled phase, the robot placed at $\bm q_k$ with joint noise and a
height jitter along the wall, and the pad latches seeded from $\bm\sigma_k$.
For a tilted wall the reference is rotated about the wall base line; joint
angles are frame-invariant.

Imitation rewards use the exponential kernel $r=\exp(-e/\sigma^2)$ on the
mean squared joint-angle and joint-velocity errors, on the root-frame
end-effector positions, and on the root-orientation geodesic error; each pose
reward is multiplied by a schedule gate
$g=g_0+(1-g_0)\bigl(1-\overline{\mathrm{mismatch}}\bigr)$, where a pad's
mismatch $|\sigma_h-\hat\sigma_h|$ is discounted linearly within a grace
fraction of its scheduled window around every edge. Schedule mismatch is
also billed directly, landings are rewarded once per window by their
distance to $\bm a_h$, and climb progress is rewarded relative to the clip's
own rate $h f_s/C$. The base's world position is never tracked: the
reference dictates \emph{shape and schedule}, the policy learns force
allocation and balance under domain-randomized pad capacities.

\section{Properties and ablations}
\label{sec:ablation}

The construction guarantees, per frame and by construction rather than by
tuning: (i)~latched pads are stationary within $\varepsilon_p,\varepsilon_r$;
(ii)~joints and gimbals stay strictly inside their limits; (iii)~released
pads and toes do not penetrate the wall; (iv)~a moment-free static force
allocation inside the limit surface and torque limits exists; (v)~dwell
times cover the engage pipeline and the recruitment ramp; (vi)~the cycle
closes exactly. We report the following ablations \tbd:

\begin{itemize}[leftmargin=1.6em]
  \item \textbf{Weighted vs.\ lexicographic IK} (Sec.~\ref{sec:weighted} vs.\
        \ref{sec:ik}): maximum latched-pad drift, worst joint-limit margin,
        number of frames reported infeasible, and the reachable envelope of
        the choreography (highest feasible hand anchor at fixed standoff).
  \item \textbf{Locked vs.\ bounded gimbal} ($\gamma_{\max}\in\{0,1,2,2.5\}^\circ$):
        pad drift and infeasible frames on the mount prologue, where the
        closed chain of two welded 4-DoF arms is over-constrained.
  \item \textbf{Physics gates on/off}: policies trained on the \texttt{fast}
        (timing-violating) and \texttt{paced} cycles; retargeted style clips
        as a further baseline.
  \item \textbf{Downstream}: engage match, per-pad occupancy, climb rate and
        fall rate of the trained policy on the vertical wall.
\end{itemize}

\end{document}
